\subsection{Universal realization and correction search over
\texorpdfstring{$\GF{13}$}{GF(13)}}

The first application starts from an optimal code and asks whether the
universal form exposes a useful parent.  The answer is exact: a known MDS
code has a rank-one realization, its two-coordinate reduction is self-dual,
and only one projective correction class over that parent can restore the MDS
distance.

\begin{prop}
\label{prop:universal-mds-137}
\rm
Let \(K=\GF{13}\), let \(c=5\), and let
\[
G_{14}=\left[I_7\ \middle|\
\operatorname{circ}(2,9,10,9,2,1,1)\right].
\]
This is the pure double-circulant self-dual MDS \([14,7,8]\) code recorded
in~\cite{BetsumiyaEtAl}. Put the columns of \(G_{14}\) in the order
\[
\sigma=\AppThirteenFourteenCoordinateOrder
\]
and group consecutive coordinates into seven two-coordinate blocks. A
row-equivalent universal generator is obtained from
\begin{equation}
P=\AppThirteenFourteenP,
\qquad
H=\AppThirteenFourteenH,
\qquad
Q=\AppThirteenFourteenQ.
\label{eq:gf13-14-coefficients}
\end{equation}
Let \(\delta_{ij}\) denote the Kronecker delta.  In the notation of
Theorem~\ref{thm:universal-rank-boxed},
\begin{equation}
\widehat G_{14}=G(5;P,H,Q,(1))
=\left[
\begin{array}{c|c}
\bigl((P_{ij},5P_{ij}+\delta_{ij})\bigr)_{1\le i,j\le6}
  & \bigl((H_i,5H_i+Q_i)\bigr)_{1\le i\le6}\\
\hline
\bigl(-Q_j(1,5)\bigr)_{1\le j\le6} & (1,5)
\end{array}
\right].
\label{eq:gf13-14-structured-generator}
\end{equation}
It has \(k=6\), \(r=1\), and \(D=(1)\).  Deleting its first block row and
first block column gives a self-dual \([12,6,6]\) parent.  For this parent,
the projective correction classes satisfying the necessary weight-six/seven
lifting condition form the singleton
\[
\Lambda_8=\{\AppThirteenFourteenSurvivor\}.
\]
The correction column of \(\widehat G_{14}\) is
\(\gamma=\AppThirteenFourteenGamma\), whose projective class is this unique
element; rebuilding recovers \(\widehat G_{14}\).
\end{prop}

\begin{proof}
Direct calculation gives \(G_{14}G_{14}^{\mathsf T}=0\) and
\(\rank G_{14}=7\). Every one of its
\(\binom{14}{7}=3432\) seven-column submatrices has rank~7. Hence every
seven columns are an information set, so the generated code is MDS and has
parameters \([14,7,8]\).  In particular, its minimum-weight coefficient is
the MDS value
\[
A_8=(13-1)\binom{14}{8}=36036.
\]

For the asserted row equivalence, if \(G_{14}^{\sigma}\) denotes the stated
column reordering, then
\[
T=\begin{pmatrix}
12&8&12&5&8&10&0\\
10&11&11&4&3&6&0\\
4&7&2&3&10&9&0\\
9&7&0&9&0&0&0\\
0&0&9&0&7&9&0\\
4&7&0&9&0&5&0\\
0&0&0&0&0&10&7
\end{pmatrix},
\qquad \det T=5,
\]
and exact multiplication gives
\(TG_{14}^{\sigma}=\widehat G_{14}\), with \(\widehat G_{14}\) as in
\eqref{eq:gf13-14-structured-generator}. Applying
\(\partial(s,t)=t-5s\) blockwise gives
\[
\partial\widehat G_{14}=\AppThirteenFourteenReadout.
\]
Thus \(\rank(\partial\widehat G_{14})=6\), while its last row lies in
\(U_5\) and ends in \((1,5)\); hence \(r=1\) and \(D=(1)\).
The universal Gram relation holds for \(P,H,Q\), so
Theorem~\ref{thm:universal-rank-boxed} gives self-duality both before and
after the simultaneous row-and-block deletion.

The parent has \(A_j=0\) for \(1\le j<6\), \(A_6=960\), and
\(A_7=3744\). After projectivizing its weight-six and weight-seven
coefficient vectors, there are \(\AppThirteenFourteenLowLines\) classes.
Exact enumeration of \(\mathbf P^5(\F_{13})\) shows that precisely
\(\AppThirteenFourteenSurvivor\) is nonorthogonal to every one of them.
Finally,
\[
[\AppThirteenFourteenGamma]=\AppThirteenFourteenSurvivor,
\]
and the lower six rows of \eqref{eq:gf13-14-structured-generator}, after
deleting their first block, are literally the parent generator. Reattaching
that block and the first row therefore recovers \(\widehat G_{14}\).
\end{proof}

This calculation supplies a search principle rather than a claim that the
known MDS code itself is new.  Starting from a fixed parent, low-weight
coefficient lines eliminate projective corrections before full minimum
distance is computed.  Here the filter is sharp: exactly one projective class
survives and it reconstructs the optimal child.  For \(m\ge0\), one may
therefore study
\[
\delta_m(G_{14})=\max\{d(C_m): C_{i+1}\text{ is obtained from }C_i
\text{ by a non-direct-sum step of
Lemma~\ref{lem:rank-boxed-building-up}}\}.
\]
The question is whether this fixed-seed search reaches the best possible
minimum distance at length \(14+2m\).  The comparison with known parameters
and the successive constructions below separate the proved examples from this forward search
question.

\FloatBarrier
